\section*{Introduction}
\addcontentsline{toc}{section}{Introduction}
The three charged leptons have identical electric charge and spin but sharply different masses. Their mass ratios are accurately reported, yet the origin of the hierarchy is not fixed by those reports themselves. Contemporary particle data organize the electron, muon, and tau as a three-member charged-lepton family \cite{PDG2024}, while precision constants provide conventional comparison standards for their masses and ratios \cite{CODATA2022}. The question addressed here is whether a definite three-rank mass structure can be obtained from a more primitive relational construction without using the measured muon or tau mass as an algebraic input.

The strategy is closure-first. The paper begins with relational distinguishability and a positive count-and-transfer discipline. Persistence, closure, and their bookkeeping consequences are developed before geometric, observer, or empirical language is admitted. A face-centered-cubic presentation with rhombic-dodecahedral cells is then obtained as a downstream relational realization. The matrices used in this development are organizational displays: independently established quantities are placed into typed rows and columns after derivation. Their rectangular arrangement does not generate a scientific relation, and no matrix algebra is used to obtain the charged-lepton result.

Within the resulting geometry, one completed cell is designated as the observer/reference role. \textbf{Observer-rooted asymmetry is the dependence of a lawful report on the observer's rooted ancestry and betweenness position, while the underlying unordered intrinsic spectrum remains presentation invariant.} The observer therefore does not change an intrinsic mass or select a preferred spectral value. Its rooted occurrence fixes the lawful presentation data, and its unique interior ancestry position determines which already formed rank receives the later finite-screen rendering.

The principal mathematical contributions are a single connected derivational sequence. A rooted two-active/one-idle occurrence fixes both a normalized radius and an action aperture. The action supplies the phase of six presentation-equivalent Hermitian operators and hence one intrinsic unordered three-amplitude spectrum. A construction-derived adjoint forces quadratic support with coefficient \(m=1/(2\pi)\). Exact rational bounds then place the middle-to-minimum intrinsic ratio strictly between \(206\) and \(207\). A recoverable one-event/two-preserved-face record supports two \(206\)-element sets sharing one three-root carrier; their free pushout has \(409\) elements. The corresponding regular-screen factor acts only on the rank selected by observer-rooted betweenness. These steps determine the two dimensionless charged-lepton mass ratios. Figure \(\ref{fig:derivation-architecture}\) records this dependency order without adding an inference.

The distinction between derivation and calibration is strict. The spectrum, rank ordering, quadratic-support law, \(206\), \(409\), finite-screen factor, and both dimensionless ratios are derived before any dimensional standard is introduced. The electron supplies the sole conventional reporting scale. Neither a measured muon value nor a measured tau value appears as an algebraic operand. The resulting muon and tau mass-energy equivalents are therefore downstream outputs relative to the electron standard, not independent absolute-mass derivations.

The development history also matters for claim strength. The \(409\) construction was developed after exposure to the muon result. The paper does not present the middle output as a blind held-out prediction, even though the measured muon value is absent from the algebra and no continuous numerical parameter is fitted. A post-derivation comparison with the CODATA recommended adjusted muon-to-electron ratio is reported transparently on that basis.

Koide's charged-lepton relation \cite{KOIDE1983} is likewise not a construction input. It is evaluated only after the mass-ratio report has closed. The word \emph{independent} in the subtitle therefore denotes logical independence from the construction inputs: Koide's relation is not used to choose the phase, radius, rank assignment, integer counts, finite-screen factor, or reporting scale. It does not assert algebraic independence from the derived radius and quadratic-support law.

The paper is organized as follows. Sections 1 and 2 establish relational admission and positive transfer arithmetic. Sections 3--7 develop closure, bookkeeping, geometry, and registration while keeping presentation distinct from authority. Sections 8--11 establish the observer, readout, and one-event/two-face recovery record. Section 12 derives the intrinsic three-rank spectrum and quadratic-support law. Section 13 proves the exact \(206\) interval, constructs \(409\), and establishes observer-rooted middle placement. Section 14 forms the charged-lepton ratios and attaches the electron reporting standard before making the disclosed post-exposure muon comparison. Section 15 evaluates the exact intrinsic and finite-screened Koide quotients as downstream consistency checks.

